\documentclass[11pt, a4paper]{scrartcl}
\usepackage{xcolor}
\usepackage[colorlinks=true,allcolors=blue,pdfborder={0 0 0},pdfstartview={FitV},breaklinks,linktocpage]{hyperref}
\usepackage{color,amsmath,amsfonts,amssymb}
\usepackage[utf8]{inputenc}
\usepackage{enumitem}%für mehrspaltige Aufzählungen
\usepackage[ngerman, english]{babel}
\usepackage{epsfig}
\usepackage{pstricks}
\usepackage{graphics}
\graphicspath{{Figures/}}
\usepackage{bbm}
%\usepackage{showlabels}
\usepackage{siunitx}
\usepackage[labelfont=bf]{caption}
%\usepackage{nonfloat}
\usepackage{booktabs}%für Tabellen
\usepackage{pgffor} %for loops
\usepackage{ifthen} %for if
\usepackage{pgfplots} %für Plots
\pgfplotsset{
compat=1.5,
every axis/.append style={line width=1.0pt, tick style={line width=1.0pt}}
} %für richtigen Abstand der labels
\usetikzlibrary{shapes,arrows,fit,calc,positioning,arrows}
\usepackage[
  textheight=25cm,
left=2.3cm,
right=2.3cm,
headheight=25pt,
  includehead,includefoot,
  heightrounded,
]{geometry}
\makeatletter
\def\input@path{{./Figures/}}
\usepackage{fancyhdr}
\pagestyle{fancy}
\fancyfoot{}
\fancyhead{}
\fancyhead[L]{}
\fancyhead[R]{\thepage}
\fancyhead[C]{\nouppercase{\leftmark}}
\usepackage{url}
\usepackage{amsthm}
\usepackage{graphicx}
\usepackage{tikzpagenodes}

\definecolor{darkgreen}{rgb}{0.0,0.6,0.0}
\definecolor{shadedgreen}{rgb}{0.8,0.95,0.8}
\definecolor{shadedblue}{rgb}{0.85,0.85,1.0}
\definecolor{darkblue}{rgb}{0.0,0.0,0.5}


\newlength\figureheight 
\newlength\figurewidth


\usepackage[varg]{txfonts} %font/Schrift





\definecolor{darkgreen}{rgb}{0.0,0.5,0.0}

% Aliases
%%%%%%%%%%%%%%%%%%%%%%%%%%%%%%%%%%%%%%%%%%%%%%%%%%%%%%%%%%%%%%%%%%
%% The Aliases defined here are used in the LaTeX template files
%% AND the python options, so that they can be used in python plot
%% titles and descriptions. See ../latex_settings_units.py
%% Lines starting with '%' are ignored by the parser.
%%%%%%%%%%%%%%%%%%%%%%%%%%%%%%%%%%%%%%%%%%%%%%%%%%%%%%%%%%%%%%%%%%
%%%%%%%%%%%%%%%%%%%%%%%%%%%%%%
%% New Operators in math mode
%%%%%%%%%%%%%%%%%%%%%%%%%%%%%%
\DeclareMathOperator{\rRe}{Re}
\DeclareMathOperator{\rIm}{Im}
%%%%%%%%%%%%%%%%%%%%%%%%%%%%%%
%$ Roman, Bold etc.
%%%%%%%%%%%%%%%%%%%%%%%%%%%%%%
\newcommand{\mc}[1]{\mathcal{#1}}
\newcommand{\mf}[1]{\mathfrak{#1}}
\newcommand{\mr}[1]{\mathrm{#1}}
\newcommand{\mb}[1]{\mathbf{#1}}
%%%%%%%%%%%%%%%%%%%%%%%%%%%%%%
%% Math symbols
%%%%%%%%%%%%%%%%%%%%%%%%%%%%%%
\newcommand{\bE}{\mb{E}}
\newcommand{\bk}{\mb{k}}
\newcommand{\bA}{\mb{A}}
\newcommand{\bj}{\mb{j}}
\newcommand{\bd}{\mb{d}}
\newcommand{\eps}{\epsilon}
\newcommand{\w}{\omega}
\newcommand{\sd}{\hspace{0.05em}}
\newcommand{\eqt}{\,{=}\,}
\newcommand{\pt}{\,{+}\,}
\newcommand{\coloneqqt}{\,{\coloneqq}\,}
\newcommand{\intbzdkpi}{\int\limits_\text{BZ}\frac{d\bk}{(2\pi)^2}}
\newcommand{\rhonnprime}{\rho_{nn'}(\bk,t)}
\newcommand{\rhonprimen}{\rho_{n'n}(\bk,t)}
\newcommand{\un}{\underline{n}}
\newcommand{\braunk}{\langle u_{n\bk}|}
\newcommand{\ketunk}{|u_{n\bk}\rangle}
\newcommand{\ketunkprime}{|u_{n'\bk}\rangle}
\newcommand{\timest}{\,{\times}\,}
%%%%%%%%%%%%%%%%%%%%%%%%%%%%%%
%% Units (LaTeX - siunitx)
%%%%%%%%%%%%%%%%%%%%%%%%%%%%%%
\newcommand{\fs}{\femto\seconds}
\newcommand{\As}{\angstrom}
\newcommand{\eV}{\electronvolt}
\newcommand{\e}{\elementarycharge}
\newcommand{\THz}{\tera\hertz}
\newcommand{\MV}{\mega\volt}

\newcommand{\paper}[4]{\item #1, \,\textit{#2}, \,\href{#3}{#4}.\\[-1.4em]}
\newcommand{\GW}{\textit{GW} }

\newcommand{\XMAXBSSCF}{1.3434623097571412}
\newcommand{\YMINBSSCF}{-3.5}
\newcommand{\YMAXBSSCF}{5.33304264}
\newcommand{\PLOTSBSSCF}{
\addplot[very thick, darkblue, smooth] table {data/band_SCF_7.dat};
\addplot[very thick, darkblue, smooth] table {data/band_SCF_8.dat};
\addplot[very thick, darkblue, smooth] table {data/band_SCF_9.dat};
\addplot[very thick, darkblue, smooth] table {data/band_SCF_10.dat};
\addplot[very thick, darkblue, smooth] table {data/band_SCF_11.dat};
\addplot[very thick, darkblue, smooth] table {data/band_SCF_12.dat};
\addplot[very thick, darkblue, smooth] table {data/band_SCF_13.dat};
\addplot[very thick, darkblue, smooth] table {data/band_SCF_14.dat};
\addplot[very thick, darkblue, smooth] table {data/band_SCF_15.dat};
\addplot[very thick, darkblue, smooth] table {data/band_SCF_16.dat};
}
\newcommand{\XTICKLABELS}{xticklabels={$\Gamma$, K, M, $\Gamma$},}
\newcommand{\TICKSSPECIALKPOINTSBSSCF}{0.0, 0.4709331162702407, 0.8434623097571405, 1.3434623097571405}


\begin{document}

\pagenumbering{Roman}

\begin{titlepage}\pdfbookmark[0]{Title}{Title}
  \sffamily
  \begin{center}
{
\includegraphics[width=5cm]{cp2k_logo.png}
\\[1em]
\Huge \bfseries Output summary of a \GW calculation\\[0.4em] using the CP2K package}
  \end{center}{\large
  \vspace{3em}
  Contact:
  \\[1em]
    Jan Wilhelm
      \\[0.5em]
    Institute of Theoretical Physics
    \\[0.5em]
    University of Regensburg
    \\[0.5em]
    jan.wilhelm@ur.de
    \\[3em]
    
\pdfbookmark[0]{Contents}{Contents}
\tableofcontents
  }
\end{titlepage}







\pagenumbering{arabic}
\pagestyle{plain}






\pagestyle{fancy}

\section{Input parameters}


\section{Band structure}
\begin{figure}[h!]
\centering
\setlength\figureheight{12cm} 
\setlength\figurewidth{\textwidth}
\begin{tikzpicture}
\begin{axis}[
height=\figureheight,
width=\figurewidth,
tick align=outside,
tick pos=left,
x grid style={white!70!black},
xmajorgrids,
xmin=0, 
xmax=\XMAXBSSCF,
xtick style={color=black},
xtick={\TICKSSPECIALKPOINTSBSSCF},
\XTICKLABELS
%xticklabels={$\Gamma$,K,$\Gamma$},
y grid style={white!70!black},
ylabel={$E-E_\text{VBM}$ in eV},
ymajorgrids,
ymin=\YMINBSSCF,
ymax=\YMAXBSSCF,
ytick style={color=black}
]
\PLOTSBSSCF
\end{axis}
\end{tikzpicture}

\caption{Electronic band structure from the self-consistent field (SCF) electronic structure method. The energy of the bands is given relative to the valence band maximum (VBM).}
    \label{fig:Efield}
\end{figure}




\section{How to reference the \GW calculation}
When using the \GW implementation in CP2K, please cite the following publications:
\begin{enumerate}[leftmargin=*]

\item[[1{]}] M.~Graml, K.~Zollner, D.~Hernangómez-Pérez, P.~E.~Faria Junior, and J.~Wilhelm, \textit{Low-scaling GW algorithm applied to twisted transition-metal dichalcogenide heterobilayers}, \href{
https://doi.org/10.48550/arXiv.2306.16066}{arXiv 2306.16066 (2023)}.

\item[[2{]}] J.~Wilhelm, P.~Seewald, D.~Golze, \textit{Low-Scaling GW with Benchmark Accuracy and Application to Phosphorene Nanosheets}, \href{https://doi.org/10.1021/acs.jctc.0c01282}{
J.~Chem.~Theory Comput.~\textbf{17}, 1662-1677 (2021)}.


\item[[3{]}] T.\;D.\;Kühne, M.\;Iannuzzi, M.\;Del\;Ben, V.\;V.\;Rybkin, P.\;Seewald, F.\;Stein, T.\;Laino, R.\;Z. Khaliullin, O.\;Schütt, F.\;Schiffmann, D.\;Golze, J.\;Wilhelm, S.\;Chulkov, M.\,H.\,Bani-Hashe\-mian, V.\,Weber, \linebreak U.\,Borst\-nik, M.\,Taillefumier, A.\,S.\,Jakobovits, A.\,Lazzaro, H.\,Pabst,  T.\,Müller,  R.\,Schade, M.\,Guidon, S.\,Ander\-matt, N.\,Holmberg, G.\,K.\,Schenter, A.\,Hehn, A.\,Bussy,  F.\,Belleflamme, G.\,Tabacchi, A.\,Glöß, M.\,Lass, I.\,Bethune, C.\,J.\,Mundy, C.\,Plessl, M.\,Watkins, J.\,VandeVondele, M.\,Krack, J.\,Hutter, \textit{CP2K: An electronic structure and molecular dynamics software package - Quickstep: Efficient and accurate electronic structure calculations}, \href{https://doi.org/10.1063/5.0007045}{J.~Chem.~Phys.~\textbf{152}, 194103 (2020)}.


\end{enumerate}

\end{document}
